\documentclass[11pt]{article}

\usepackage[margin=1in]{geometry}
\usepackage{amsmath,amssymb,amsthm}
\usepackage{booktabs}
\usepackage{microtype}
\usepackage[T1]{fontenc}
\IfFileExists{newtxtext.sty}{\usepackage{newtxtext}\usepackage{newtxmath}}{\usepackage{lmodern}}
\usepackage{natbib}
\usepackage[hidelinks]{hyperref}

\newcommand{\Fp}{\mathcal F_{t-1}}
\newcommand{\Ee}{\mathbb E}

\title{\vspace{-2em}Spectral Calibration Without a Split\\[0.3em]
       \large Learning Covariance Eigenvalue Corrections One Observation at a Time}
\author{Peter Cotton}
\date{Note --- \texttt{microprediction/precise}}

\begin{document}
\maketitle

\begin{abstract}
\noindent
Cross-validated covariance cleaning learns the variance belonging to an empirical eigenvector by
projecting held-out data onto it, which requires partitioning a sample. A stream needs no partition.
The eigenvector basis carried forward from past data is already out of sample with respect to the
observation about to arrive, so every observation labels every eigenvalue at once. We give the
conditional identity this rests on and the one-pass estimator that follows, which costs $O(p^2)$ per
observation and keeps bounded state. In five regimes the learned calibration beats an uncleaned
exponentially weighted covariance everywhere, and a simpler estimator beats it wherever dependence
drifts: the analytical map applied to a rolling window, where equal weights make the asymptotics
exact with $n = W$.
\end{abstract}

\section{The gap}

Nonlinear shrinkage of a covariance spectrum, which moves each sample eigenvalue by its own amount
while keeping the sample eigenvectors, has an analytical solution under equally weighted, identically
distributed observations \citep{ledoitpeche2011,ledoitwolf2020analytical}, and a random-matrix
counterpart in the rotationally invariant estimators \citep{bun2017cleaning}. Both derivations assume
a fixed population covariance observed through sampling noise. Financial dependence is not fixed. It
drifts, and common modes whose strength varies dominate it. \citet{manolakis2026} make this
the explicit motivation for replacing the analytical singular-value map with a learned one, fixing
the empirical basis from theory so that only the part where the derivation is known to fail is
learned. Their map is trained offline, on a corpus.

The weighted case has since been worked out analytically: \citet{oriol2024weighted} give non-linear
shrinkage formulas for weighted sample covariances, the exponentially weighted case explicitly. That
is the principled route for a recency-weighted analytical cleaner, and we would take it. We ask
instead whether the correction can be \emph{learned from the stream itself}, with bounded state, no
corpus and no pretraining.

\section{The identity}

Let $\Fp$ denote everything observable before time $t$, and let $U_{t-1} = [u_1 \cdots u_p]$ be an
orthonormal basis computed entirely from $\Fp$. Suppose $\Ee[x_t \mid \Fp] = 0$ and
$\operatorname{Cov}(x_t \mid \Fp) = \Sigma_t$. Project the new observation,
\[
  z_t = U_{t-1}^\top x_t .
\]
Because $u_i$ is $\Fp$-measurable it passes through the conditional expectation, giving
\begin{equation}\label{eq:label}
  \Ee\!\left[z_{i,t}^2 \mid \Fp\right] \;=\; u_i^\top \Sigma_t\, u_i ,
\end{equation}
which is exactly the variance a cleaned spectrum must assign to that empirical eigenvector. One
observation supplies a noisy but conditionally unbiased label for \emph{every} eigenvalue
simultaneously. Equation \eqref{eq:label} requires neither Gaussianity nor stationarity; it is a
statement about conditioning, and the only structural demand is that $U_{t-1}$ not be built using
$x_t$.

That the squared projections are the right thing to predict follows from an exact decomposition. For
any orthogonal $U$, any $d \in \mathbb R^p$, and $z = U^\top x$,
\begin{equation}\label{eq:frob}
  \bigl\| U \operatorname{diag}(d) U^\top - x x^\top \bigr\|_F^2
   \;=\; \sum_i \bigl(d_i - z_i^2\bigr)^2 \;+\; \underbrace{\sum_{i \neq j} z_i^2 z_j^2}_{\text{free of } d} .
\end{equation}
The Frobenius norm is invariant under orthogonal conjugation, and
$U^\top ( U \operatorname{diag}(d) U^\top - xx^\top ) U = \operatorname{diag}(d) - z z^\top$, whose
diagonal entries are $d_i - z_i^2$ and whose off-diagonal entries are $-z_i z_j$. So learning to
predict squared projections minimizes, term for term, the same instantaneous objective as learning
the covariance matrix within the chosen basis. The residual term is the price of holding the
eigenvectors fixed, and no choice of $d$ can touch it.

Neither statement is new. Equation \eqref{eq:label} is what makes cross-validated cleaning work:
\citet{bartz2016} estimates cleaned eigenvalues by projecting held-out folds onto eigenvectors fitted
on the remainder, and \citet{lamrani2025holdout} analyze how such a sample should optimally be split,
finding the train--test ratio should scale as the square root of the dimension. The observation we
have not found recorded is that in a streaming setting the split need not be constructed at all.
Every arriving observation is already held out with respect to the basis carried forward, so the
data cost of cross-validation, the reason one must decide how much sample to sacrifice, simply does
not arise. What was a design choice becomes a property of the arrow of time.

\section{The estimator}

State: a running mean $m$, an exponentially weighted second moment $S$, a basis $U$, a spectrum
$\lambda$, and the accumulators of a calibration map $f$. On observing $x_t$, with weight $w$:
\begin{align*}
  \delta &= x_t - m, \qquad z = U^\top \delta, \\
  &\text{train } f \text{ on the pairs } (\text{features}_i,\; z_i^2), \\
  \lambda &\leftarrow (1-w)\lambda + w z^2, \\
  m &\leftarrow m + w\delta, \qquad S \leftarrow (1-w) S + w\, \delta \delta^\top,
\end{align*}
refreshing $(\lambda, U)$ from an eigendecomposition of $S$ every $K$ observations, and reporting
$\hat\Sigma_t = U \operatorname{diag}\bigl(f(\lambda)\bigr) U^\top$. The ordering matters: the labels
are formed against the basis and the mean that predate $x_t$, which is what makes \eqref{eq:label}
apply.

The $\lambda$ recursion is not an approximation. For a basis held fixed between refreshes,
\[
  u_i^\top S_t u_i = (1-w)\, u_i^\top S_{t-1} u_i + w \bigl(u_i^\top \delta\bigr)^2 ,
\]
so the exponentially weighted Rayleigh quotient obeys the same recursion as the covariance and is
available at $O(p)$ from projections already computed. Only the basis needs an eigendecomposition.
The per-observation cost is therefore $O(p^2)$ for the projection and the outer-product
update, plus $O(p^3/K)$ amortized, and nothing grows with elapsed stream length. Stale eigenvectors
are the price. They bound how quickly the estimator can follow a rotation of the true eigenbasis.

\section{What has to be learned}

The map $f$ must be positive, must treat equal eigenvalues equally, and must forget. We use
kernel-weighted running averages over a fixed grid of feature centers, a Nadaraya--Watson estimator
with exponential forgetting, so that the prediction is a ratio of two nonnegative accumulators. It
is positive because the labels are squares, it pools directions with identical features by
construction, and it forgets at a controlled rate. A few dozen bins suffice; nothing here needs a
neural framework.

The design question is what to index the map by. Three choices separate cleanly:
\emph{rank}, the position of an eigenvalue in the sorted spectrum; \emph{eigenvalue}, its observed
size relative to the mean, on a log scale; and \emph{shared}, the observed size together with a
global summary of the spectrum's shape (we use the dispersion of the log spectrum). Rank calibration
can memorize a useful spectral profile; the concern is that it will keep imposing that profile after
the structure generating it is gone, since the position of an eigenvalue says nothing about whether
it still refers to anything.

\section{Experiments}

Gaussian streams in $p = 32$ dimensions, $6000$ observations, eight seeds, exponential weight
$r = 0.02$, basis refreshed every $10$ observations, calibration forgetting $\rho = 0.02$, scoring
from observation $1000$ onward. The metric is mean squared Frobenius error divided by the squared
norm of the truth; lower is better. Five regimes: an isotropic covariance; a broad, static
population spectrum; an abrupt rotation of the factor directions at the midpoint; a factor structure
that disappears at the midpoint; and one in which the factor structure and its absence alternate
repeatedly, the only regime in which a spectral shape \emph{recurs} and can therefore be recalled.

We compare against an uncleaned exponentially weighted covariance, the linear shrinkers
\texttt{LedoitWolf} and \texttt{OAS} at the same decay, and the analytical nonlinear shrinker of
\citet{ledoitwolf2020analytical} over an expanding sample.

\begin{table}[ht]
\centering
\small
\begin{tabular}{lrrrrrrrr}
\toprule
 & uncleaned & \multicolumn{2}{c}{linear, weighted} & \multicolumn{2}{c}{analytical nonlinear} & \multicolumn{3}{c}{learned calibration} \\
\cmidrule(lr){2-2}\cmidrule(lr){3-4}\cmidrule(lr){5-6}\cmidrule(lr){7-9}
regime & raw EW & L--W & OAS & expanding & window & rank & eigenvalue & shared \\
\midrule
isotropic         & 0.340 & \textbf{0.0007} & \textbf{0.0007} & 0.0007          & 0.0049          & 0.0026 & 0.0018 & 0.0018 \\
broad spectrum    & 0.206 & 0.204           & 0.199           & \textbf{0.0072} & 0.066           & 0.135  & 0.147  & 0.147  \\
factor rotation   & 0.243 & 0.238           & 0.231           & 0.206           & \textbf{0.067}  & 0.191  & 0.186  & 0.186  \\
factors vanish    & 0.299 & 0.127           & 0.124           & 0.239           & \textbf{0.059}  & 0.105  & 0.103  & 0.103  \\
regimes alternate & 0.330 & 0.131           & 0.128           & 0.217           & \textbf{0.104}  & 0.114  & 0.109  & 0.109  \\
\bottomrule
\end{tabular}
\caption{Mean squared Frobenius error relative to the squared norm of the truth; lower is better,
best in each row in bold. Eight seeds, $p=32$, $6000$ observations each. The windowed
column uses $W=250$; the learned columns use the same exponential weight $r=0.02$ as the linear
shrinkers.}
\end{table}

\emph{The learned correction is dominated by a simpler analytical one.} On every regime where the
dependence changes, applying the analytical map to a rolling window beats every learned variant:
$0.067$ against $0.186$ under a rotation, $0.059$ against $0.103$ when the factor structure dies,
$0.104$ against $0.109$ under alternation.

The learned calibration does what we set out to show it could. It beats an uncleaned exponentially
weighted covariance everywhere, and beats exponentially weighted linear shrinkage by fifteen to
twenty percent on exactly the non-stationary regimes that motivated it. But linear shrinkage was the
wrong bar. A rolling window gives the analytical nonlinear map a sample it can forget without
approximating anything, since equal weights inside a window of length $W$ are exactly the sample the
asymptotics describe, and that cheap route captures the adaptivity the learned route was introduced
to supply.

The gap narrows to five percent on the fastest-drifting regime. That is the one direction in which a
case for learning could be rebuilt: the window's advantage should shrink as the drift rate approaches
the window length, and a learned map need not commit to a timescale in advance.

\emph{Rank calibration remembers, and that cuts both ways.} Indexing by spectral position is the
best learned variant on the static broad spectrum ($0.135$ against $0.147$), where a stable profile
pays to memorize, and the worst on every regime where the structure changes. Conditioning on the
observed eigenvalue instead removes that failure at a modest cost elsewhere, That is what one would
hope. An eigenvalue's size still refers to something after the structure that produced it is gone,
whereas its rank does not.

\emph{Global shape conditioning does not help.} We expected the third variant, indexing the map by
the observed eigenvalue \emph{and} a summary of the spectrum's overall shape, to dominate the
second, on the reasoning that it could recall a calibration when a shape recurred. It does not. The
two agree to three decimal places in every regime, including the alternating one built specifically
to reward recall. The reason is visible in the diagnostics. The shape statistic ranges over roughly
$0.3$--$1.1$ \emph{across} these regimes but moves only about $0.4$ \emph{within} any single stream,
so the discriminating variation is between problems rather than along a trajectory, and a map that
one estimator carries through one stream never sees enough of the axis to use it. We report this as
a negative result rather than tuning until it disappears.

\section{Relation to prior work}

The analytical line \citep{ledoitpeche2011,ledoitwolf2020analytical,bun2017cleaning} derives the map
rather than learning it, and is the right default whenever its assumptions hold; the expanding-sample
column above shows how much it wins by when they do, and how badly it fails when they do not. The
cross-validation line \citep{bartz2016,lamrani2025holdout} uses the same identity as we do but must
construct its split. \citet{oriol2024weighted} extends the analytical treatment to weighted samples
and is the principled way to build a recency-weighted analytical cleaner; substituting a
moment-matched effective sample size into an equally weighted formula, as is sometimes done, is a
scalar approximation that does not encode the weight distribution. \citet{manolakis2026} learn the
singular-value map, as we do, but offline and for the rectangular cross-covariance problem. For that
problem the corresponding sequential label is immediate,
$(u_i^\top x_t)(v_i^\top y_t)$, though assembling a valid joint covariance additionally requires
compatibility with the marginal blocks: an arbitrary cross-block correction need not preserve
positive semidefiniteness.

Incremental eigenspace tracking is a separate and well-developed literature
\citep{weng2003ccipca}; we take the basis by periodic refresh, which is crude, and a better tracker
would directly relax the one cost we identified.

\section{Limitations}

The labels are severely noisy. Under Gaussianity $z_i^2$ has variance $2(u_i^\top \Sigma u_i)^2$, so
a single observation carries roughly one effective bit per direction, and the map's forgetting rate
trades adaptation against that noise. The eigenvector basis is refreshed rather than tracked, which
bounds adaptation to abrupt rotation.

All evidence here is synthetic and Gaussian, and we constructed the regimes to separate the failure
modes we suspected. There is no real equity study. On those regimes the construction does not win: a
rolling window carrying the analytical map beats it wherever the dependence drifts, which is the case
the learned map was meant to answer.

What survives is the identity, the fact that a stream needs no split to exploit it, and an estimator
cheap enough that the question can be revisited on real data at drift rates faster than a window can
follow. The open question is not whether spectral cleaning can be made online. It plainly can. It is
how calibration information should be shared across directions and forgotten as the spectrum changes,
and the negative result on shape conditioning says the obvious first answer is not the right one.

\bibliographystyle{plainnat}
\bibliography{../refs}  % shared across manuscripts

\end{document}
